\chapter{Introduction}

\section{The central question}
\label{sec:intro:central}
% =============================================================================

Statistical inference is, at its core, a problem of rational learning.
Given a sequence of observations, an agent must form beliefs about the
world, update them as new data arrives, and act on them in a principled
way.  Bayesian decision theory provides the normative framework for this
process \citep{Wald1950, Savage1954, Berger1985}: beliefs are represented
as probability distributions, updating follows Bayes' rule, and actions
are chosen to minimize expected loss under the current belief state.  This
framework is both a philosophical position; that uncertainty should be
quantified probabilistically, and a practical decree for how
inference algorithms should be built.

This thesis studies three facets of Bayesian learning through the lens of
Bayesian nonparametrics.  The first concerns dynamics: how should a
Bayesian agent update its beliefs recursively, one observation at a time,
and what geometric structure underlies that recursion. The second concerns compression: when the full dataset cannot be retained, what is the
minimal summary that preserves the agent's predictive accuracy?  The third
concerns pooling: when data arrive from multiple heterogeneous
sources, how should the agent borrow strength across them without imposing
unjustified parametric constraints.  Each question has a partial answer
in the statistics literature; the implied
posterior is not identified in the first case, exact inference scales
poorly with~$n$ in the second, and parametric families impose unjustified
homogeneity assumptions in the third, and in each case Bayesian
nonparametric tools provide the language needed to sharpen and resolve the
incompleteness.

% =============================================================================
\section{Bayesian learning as a decision-theoretic framework}
\label{sec:intro:framework}
% =============================================================================

% -----------------------------------------------------------------------------
\subsection{Bayes' rule as optimal belief updating}
\label{subsec:intro:bayes}
% -----------------------------------------------------------------------------

Let $x_{1:n} = (x_1, \ldots, x_n)$ denote a sequence of observations
taking values in $\cX$, and let $G \in \cP(\Th)$ denote an unknown
probability measure on a parameter space $\Th$, indexed by a family of
likelihoods $\{k(\,\cdot\mid\theta)\}_{\theta\in\Th}$.  A Bayesian agent
begins with a prior $\pi(G)$ encoding initial beliefs about $G$ and
updates it upon observing $x_{1:n}$ to obtain the posterior
\begin{equation}
  \pi(G \mid x_{1:n})
  \;\propto\;
  \pi(G)\; p(x_{1:n} \mid G),
  \label{eq:bayes-rule}
\end{equation}
where $p(x_{1:n} \mid G) = \int \prod_{i=1}^{n} k(x_i \mid \theta)\,
G(d\theta)$ under the exchangeable i.i.d.\ assumption, which we relax in
Section~\ref{subsec:intro:definetti}.  An alternative view from the optimization persective is Donsker-Varadhan's formula \citep{donskervaradhan}. It states that the posterior is the unique belief
state minimising the expected Kullback--Leibler divergence from the true
data-generating distribution, subject to the prior as a regularizer
\citep{Zellner1988},
\begin{equation}
  \pi(\,\cdot\mid x_{1:n})
  \;=\;
  \argmin_{Q\,\in\,\cP(\cP(\Th))}
  \bE_Q\!\left[-\log p(x_{1:n}\mid G)\right]
  + \KL{Q}{\pi}.
  \label{eq:bayes-variational}
\end{equation}
This variational characterization of the posterior provides
the language for judging approximate inference algorithms. Making that
implicit objective explicit situates them from a Bayesian perspective,
and it is the lens through which we analyze Newton's recursion in
Chapter~\ref{chap:gradientflows}.

% -----------------------------------------------------------------------------
\subsection{The predictive view of Bayesian inference}
\label{subsec:intro:predictive}
% -----------------------------------------------------------------------------

The posterior $\pi(G \mid x_{1:n})$ is the conventional object of Bayesian
inference, but it is not the only complete characterization.  An equivalent
and in many ways more fundamental representation operates entirely at the
level of observables, without reference to the prior $\pi$. The
sequence of one-step-ahead predictive distributions,
\begin{equation}
  p(x_{n+1} \mid x_{1:n})
  \;=\;
  \int k(x_{n+1} \mid \theta)\;
  \pi(d\theta \mid x_{1:n}),
  \qquad n = 0, 1, 2, \ldots,
  \label{eq:predictive}
\end{equation}
completely determines the posterior under regularity conditions: given the
full sequence of predictives, one can recover $\pi(d\theta \mid x_{1:n})$ for
any $n$ \citep[Theorem~1]{FongHolmesWalker2023}.   The
Bayesian agent's beliefs about $\theta$ are entirely expressed in its
predictions about future data; the posterior contains no inferential
content beyond what the predictive sequence already encodes in the sense
that the posterior is the almost sure limit of the sequence of empirical
measures induced by the predictives as $n \to \infty$.

This predictive-first perspective decouples inference from the
specification of a parameter space $\Th$, which may be
infinite-dimensional or not explicitly given, and provides a clean
decision-theoretic criterion for evaluating approximations: an approximate
posterior is assessable by whether it produces accurate predictive
distributions, a criterion that is empirically measurable and connects
naturally to prequential scoring rules \citep{Dawid1984}.

% -----------------------------------------------------------------------------
\subsection{Martingale posteriors and posterior-free inference}
\label{subsec:intro:martingale}
% -----------------------------------------------------------------------------

The predictive-first view raises a natural question: must a valid Bayesian
analysis begin with a prior and derive predictives, or can predictives be
specified directly with the posterior as the derived object?
\citet{FongHolmesWalker2020} and \citet{WalkerHolmes2024} answer this
affirmatively.  A sequence of predictive distributions
$\{p(x_{n+1} \mid x_{1:n})\}_{n \geq 0}$ satisfying exchangeability of
the induced joint distribution over $(x_1, x_2, \ldots)$ and a martingale
condition on the sequence of induced mixing measures, namely that
$\bE[G_{n+1} \mid x_{1:n}] = G_n$ a.s., defines a valid
posterior without requiring an explicit prior specification.  The posterior
becomes a derived object: it is the almost sure limit $G_\infty =
\lim_{n\to\infty} G_n$ of the induced sequence of mixing measures, rather
than the starting point of the analysis.

This martingale posterior framework has a direct bearing on
Chapter~\ref{chap:gradientflows}.  Newton's recursion can be reinterpreted
as a specification of predictive distributions: at each step, $G_{n+1}$ is
the mixing measure that reweights the current estimate $G_n$ by the
likelihood of the new observation $x_{n+1}$.  The gradient flow analysis
in Chapter~\ref{chap:gradientflows} reveals which objective this predictive
sequence implicitly optimises, and identifies the posterior (if any(
that it targets.  As we will show, whether Newton's recursion defines a
valid martingale posterior depends critically on whether a prior
regularisation term is included in the objective.

% -----------------------------------------------------------------------------
\subsection{Sequential learning and the cost of forgetting}
\label{subsec:intro:compression}
% -----------------------------------------------------------------------------

A Bayesian agent processing a stream of observations $x_1, x_2, \ldots$
ideally retains the full posterior $\pi(G \mid x_{1:n})$ at each step.  In
parametric models this is feasible: the posterior is summaried by a
finite-dimensional sufficient statistic and storage costs are bounded.  In
nonparametric models, however, the posterior grows with $n$, the DP
posterior after $n$ observations is a random measure supported on a space
whose combinatorial complexity grows like the Bell number $B_n$, which is
super exponential in $n$ \citep{Rota1964}.  For large $n$, exact
sequential inference becomes computationally prohibitive.

This creates a fundamental tension at the heart of Bayesian nonparametrics: the expressiveness of the model, its ability to grow in
complexity with the data, is precisely what makes it expensive to
implement.  The natural question is whether the full posterior can be
replaced by a compressed summary that preserves predictive accuracy.  A
coreset is such a summary: a weighted subsample
$(x_{i_1}, w_1), \ldots, (x_{i_m}, w_m)$ of size $m \ll n$, ideally with
$m$ growing sublinearly in $n$, such that the posterior computed on the
coreset approximates the posterior on the full dataset.  Standard coreset
constructions designed for parametric models, \emph{c.f.} \citep{HugginsCampbellBroderick2016, CampbellBroderick2018}, are not well suited for non-Euclidean observations (like partitions) because they ignore the underlying topology.  Predictive matching, the approach developed in
Chapter~\ref{chap:coresets}, addresses this by defining coreset quality
directly in terms of the one-step-ahead predictive distributions,
connecting the compression problem back to the framework of
Section~\ref{subsec:intro:predictive}.

% =============================================================================
\section{Bayesian nonparametrics}
\label{sec:intro:bnp}
% =============================================================================

% -----------------------------------------------------------------------------
\subsection{Why nonparametrics?}
\label{subsec:intro:whybnp}
% -----------------------------------------------------------------------------

A parametric model assumes that the data-generating distribution belongs
to a family $\{p(\,\cdot\mid\theta) : \theta \in \mathbb{R}^d\}$ indexed
by a finite-dimensional parameter.  This is a strong assumption: it
requires the analyst to commit, before seeing the data, to a functional
form for the likelihood.  When this form is misspecified, posterior
inference can be systematically biased in ways that do not vanish as $n
\to \infty$ --- a phenomenon documented rigorously by \citet{Freedman1963}
in the context of Bayes estimators under model misspecification, and of
practical concern in density estimation, clustering, and survival
analysis.  Nonparametric models address this by letting complexity grow
with the data: the unknown object is a distribution, a density, or a
partition, and the prior is placed directly over the relevant
infinite-dimensional space.

It is worth noting that this nonparametric perspective is not confined to
the Bayesian tradition.  In machine learning, deep neural networks with
sufficiently many parameters effectively define nonparametric models over
function spaces --- in the infinite-width limit they converge to Gaussian
processes \citep{Neal1996, LeeEtAl2018}, and their implicit priors are
determined by architecture and initialisation rather than deliberate
specification.  The difference between deep learning and BNP is one of
intention: BNP chooses infinite-dimensional priors deliberately, with
known properties and principled update rules, while deep learning arrives
at nonparametric behaviour as an emergent consequence of scale.  The gap
between these two traditions --- powerful implicit priors on one side,
principled explicit priors on the other --- is one of the central open
tensions in modern statistics and machine learning.  This thesis works
squarely within the BNP tradition, though the gradient flow framework
developed in Chapter~\ref{chap:gradientflows} connects naturally to
gradient-based training of neural networks and may speak to both sides.

% -----------------------------------------------------------------------------
\subsection{De Finetti's theorem}
\label{subsec:intro:definetti}
% -----------------------------------------------------------------------------

The foundational justification for placing a prior over an
infinite-dimensional space comes not from subjective probability but from a
structural property of the observations themselves.  A sequence of random
variables $x_1, x_2, \ldots$ taking values in a Polish space $\cX$ is
exchangeable if its joint distribution is invariant under finite
permutations: $(x_1, \ldots, x_n) \eqd (x_{\sigma(1)}, \ldots,
x_{\sigma(n)})$ for all $n \geq 1$ and all permutations $\sigma$ of $[n]$.
De Finetti's theorem \citep{deFinetti1937, HewittSavage1955} states that
any infinite exchangeable sequence on a Polish space admits the
representation
\begin{equation}
  p(x_{1:n})
  \;=\;
  \int \prod_{i=1}^{n} p(x_i \mid G)\; \pi(dG),
  \label{eq:definetti}
\end{equation}
where $\pi$ is a probability measure on $\cP(\cX)$ and the product form
holds for $\pi$-almost every $G$.  The measure $\pi$ is unique and equals
the almost-sure weak limit of the empirical measures $\frac{1}{n}
\sum_{i=1}^n \delta_{x_i}$, which is a statement about the concentration
of the posterior around the truth rather than a definition of $\pi$
itself.

Exchangeability is a condition on the agent's beliefs, not on the
data-generating process< it expresses that the agent has no prior
information distinguishing one observation index from another.  Under this
condition, any coherent exchangeable
belief system must be representable as a mixture of i.i.d.\ processes, and
$\pi$ is the mixing distribution.  BNP is thus not an approximation to
some true parametric model but the natural consequence of exchangeable
uncertainty about an unknown distribution $G$.  Extensions to partial
exchangeability, where observations are exchangeable within groups but
not across them, are developed by \citet{Kingman1978} and
\citet{Aldous1985}, and are relevant to the random partition setting of
Chapter~\ref{chap:coresets}, where exchangeability holds within clusters
but the partition structure itself introduces between-cluster dependence.

% -----------------------------------------------------------------------------
\subsection{The Dirichlet process}
\label{subsec:intro:dp}
% -----------------------------------------------------------------------------

The Dirichlet process \citep{Ferguson1973} is the canonical prior over
$\cP(\cX)$ on a Polish space $\cX$.  It is parameterised by a
concentration parameter $\alpha > 0$ and a base measure $H \in \cP(\cX)$,
written $G \sim \DP{\alpha}{H}$, and defined by the requirement that for
any finite measurable partition $\{A_1, \ldots, A_k\}$ of $\cX$,
\begin{equation}
  \bigl(G(A_1), \ldots, G(A_k)\bigr)
  \;\sim\;
  \Dir\!\bigl(\alpha H(A_1), \ldots, \alpha H(A_k)\bigr).
  \label{eq:dp-def}
\end{equation}
The existence of such a process on $\cP(\cX)$ follows from Kolmogorov's
extension theorem applied to the projective system of finite-dimensional
Dirichlet distributions \citep[Theorem~1]{Ferguson1973}.  The base measure
$H$ is the prior mean --- $\bE[G(A)] = H(A)$ for all measurable $A$ ---
while $\alpha$ controls concentration around $H$: as $\alpha \to \infty$,
$G \to H$ almost surely, while small $\alpha$ allows $G$ to differ
substantially from $H$.

The DP has two properties central to its use in practice.  First, it is
conjugate to i.i.d.\ sampling: if $\theta_1, \ldots, \theta_n \mid G \iid
G$ and $G \sim \DP{\alpha}{H}$, then
\begin{equation}
  G \mid \theta_{1:n}
  \;\sim\;
  \DP{\alpha + n}{\,\dfrac{\alpha H + \sum_{i=1}^{n} \delta_{\theta_i}}{\alpha + n}},
  \label{eq:dp-posterior}
\end{equation}
where $\bigl(\alpha H + \sum_{i=1}^n \delta_{\theta_i}\bigr)/(\alpha+n)$
denotes a probability measure --- a convex combination of $H$ and the
empirical measure $n^{-1}\sum_i \delta_{\theta_i}$.  Second, realisations
of the DP are almost surely discrete: $G = \sum_{k=1}^\infty w_k
\delta_{\theta_k^*}$, where the weights follow the stick-breaking
construction
\begin{equation}
  w_k = v_k \prod_{j=1}^{k-1}(1 - v_j),
  \qquad
  v_k \iid \Beta(1, \alpha),
  \label{eq:stick-breaking}
\end{equation}
due to \citet{Sethuraman1994}.  This almost-sure discreteness induces a
natural clustering structure: observations sharing the same atom
$\theta_k^*$ are grouped into a cluster, and the predictive probability of
joining cluster $k$ is proportional to its current size --- a property
formalised by the P\'{o}lya urn scheme of
Section~\ref{subsec:intro:predictive-bnp}.

% -----------------------------------------------------------------------------
\subsection{Random partitions and species sampling models}
\label{subsec:intro:partitions}
% -----------------------------------------------------------------------------

The almost-sure discreteness of the DP established in
Section~\ref{subsec:intro:dp} has a combinatorial consequence.  When
$\theta_1, \ldots, \theta_n \mid G \iid G$ and $G \sim \DP{\alpha}{H}$,
repeated values occur with positive probability because $G$ is almost
surely discrete.  These repetitions define a random partition $\rho_n$ of
$[n] = \{1, \ldots, n\}$: indices $i$ and $j$ belong to the same block if
and only if $\theta_i = \theta_j$, which is well-defined by the almost-sure
discreteness established in Section~\ref{subsec:intro:dp}.  The first
observation always starts a new block; thereafter, given $\rho_n$ with
blocks of sizes $n_1, \ldots, n_\ell$, the $(n+1)$-th observation joins
block $k$ with probability $n_k / (\alpha + n)$ or starts a new block with
probability $\alpha / (\alpha + n)$.  This is the Chinese restaurant
process \citep{Aldous1985, Pitman2006}.

The CRP is an instance of a \emph{species sampling model} \citep{Pitman1996}:
an exchangeable random partition whose distribution is characterised by its
exchangeable partition probability function (EPPF) $\Pi_n(n_1, \ldots,
n_\ell)$, which depends on the partition only through the multiset of block
sizes $\{n_1, \ldots, n_\ell\}$, independently of which indices belong to
which block.  For the DP the EPPF takes the form
\begin{equation}
  \Pi_n(n_1, \ldots, n_\ell)
  \;=\;
  \frac{\alpha^\ell}{\alpha^{(n)}}
  \prod_{j=1}^{\ell} (n_j - 1)!,
  \label{eq:dp-eppf}
\end{equation}
where $\alpha^{(n)} = \prod_{j=0}^{n-1}(\alpha + j)$ is the rising
factorial and $(n_j - 1)! = \Gamma(n_j)$, which equals $1$ when $n_j = 1$
as required for a singleton block.  More general EPPFs, such as those of
the two-parameter Pitman--Yor process \citep{PitmanYor1997} with parameters
$(\alpha, \sigma) \in (-\sigma, \infty) \times [0, 1)$, produce
heavier-tailed cluster size distributions and arise naturally in linguistic
and genetic applications \citep{Teh2006, FavaroLijoiPrunster2012}.  The
EPPF is the central object in Chapter~\ref{chap:coresets}: the coreset
construction developed there is designed to preserve the EPPF of the full
dataset --- specifically its exchangeability structure and the recursive
product form that determines how cluster sizes grow with $n$ --- within a
weighted subsample of size $m \ll n$.

% -----------------------------------------------------------------------------
\subsection{Predictive characterisation of BNP models}
\label{subsec:intro:predictive-bnp}
% -----------------------------------------------------------------------------

Every Bayesian nonparametric model has a predictive representation, and for
many models this representation is more natural than the prior itself.  For
the Dirichlet process, the predictive distribution of $\theta_{n+1}$ given
$\theta_{1:n}$ --- obtained by marginalising $G$ out of $k(\theta_{n+1}
\mid G)$ against the posterior $\pi(G \mid \theta_{1:n})$ --- is the
P\'{o}lya urn scheme \citep{BlackwellMacQueen1973}:
\begin{equation}
  \theta_{n+1} \mid \theta_{1:n}
  \;\sim\;
  \frac{\alpha}{\alpha + n}\, H
  \;+\;
  \frac{1}{\alpha + n} \sum_{i=1}^{n} \delta_{\theta_i}.
  \label{eq:polya-urn}
\end{equation}
With probability $\alpha/(\alpha + n)$ the new observation draws a
genuinely new value from $H$, and with probability $n/(\alpha + n)$ it
repeats one of the previous values, chosen uniformly.  Applied
sequentially, this rule generates the full joint distribution of
$(\theta_1, \theta_2, \ldots)$: the partition induced by repetitions is
exactly the CRP of Section~\ref{subsec:intro:partitions}, and the posterior
update~\eqref{eq:dp-posterior} is recoverable from the predictive
weights \citep{BlackwellMacQueen1973}.

The deeper point is that this connection runs in both directions.
\citet{FortiniPetrone2012} and \citet{FongHolmesWalker2023} establish that
the predictive sequence completely determines the posterior, the prior, and
the EPPF: all inferential content of the DP is encoded in the P\'{o}lya
urn rule, and no additional information is contained in the posterior given
the full predictive sequence.  Specifying a BNP model is therefore
equivalent to specifying a coherent sequence of one-step-ahead predictive
distributions, and the posterior is the derived object obtained by taking
the almost-sure limit of the induced empirical measures.  Different BNP
models correspond to different predictive rules: the normalised
inverse-Gaussian process \citep{LijoiMenaPrunster2005}, the Pitman--Yor
process \citep{PitmanYor1997}, and the two-parameter Poisson--Dirichlet
family \citep{Pitman2006} each encode different assumptions about tail
behaviour through their predictive weights.  Newton's recursion, analysed
in Chapter~\ref{chap:gradientflows}, is a predictive rule of this type ---
the gradient flow analysis reveals which posterior, if any, its predictive
sequence targets.

% -----------------------------------------------------------------------------
\subsection{P\'{o}lya trees}
\label{subsec:intro:polytrees}
% -----------------------------------------------------------------------------

While the Dirichlet process places a prior over discrete distributions,
many inferential problems require a prior over absolutely continuous
distributions --- density estimation, goodness-of-fit testing, and the
meta-analysis of Chapter~\ref{chap:metaanalysis}, where the object of
interest is a population of smooth densities.  The P\'{o}lya tree
\citep{Lavine1992, MauldinSudderethWilliams1992} addresses this by
constructing a prior over densities through a recursive binary partitioning
of the sample space.

Formally, let $\cB = \{B_\varepsilon : \varepsilon \in \{0,1\}^*\}$ be a
nested binary partition of $\cX$, where each node $\varepsilon$ at level
$|\varepsilon|$ is split into left child $B_{\varepsilon 0}$ and right
child $B_{\varepsilon 1}$.  A P\'{o}lya tree prior $\PT{\cB}{\cA}$ assigns
to each internal node $\varepsilon$ an independent Beta random variable
$Y_\varepsilon \sim \Beta(\alpha_{\varepsilon 0}, \alpha_{\varepsilon 1})$,
where $\cA = \{\alpha_{\varepsilon j}\}$ is the collection of
hyperparameters.  The conditional probability of $B_{\varepsilon 0}$ given
$B_\varepsilon$ is $Y_\varepsilon$, and of $B_{\varepsilon 1}$ given
$B_\varepsilon$ is $1 - Y_\varepsilon$.  When the hyperparameters grow
sufficiently fast with depth --- specifically when $\alpha_{\varepsilon 0}
+ \alpha_{\varepsilon 1} \geq c \cdot |\varepsilon|^2$ for some $c > 0$
\citep{Lavine1992} --- the resulting distribution is almost surely
absolutely continuous.  Intuitively, making the Beta parameters grow with
depth forces the distribution to spread mass smoothly across the partition,
preventing concentration on a countable set of atoms, a property the DP
does not share.

P\'{o}lya trees are conjugate: observing $x_1, \ldots, x_n$ updates each
$\Beta(\alpha_{\varepsilon 0}, \alpha_{\varepsilon 1})$ to
$\Beta(\alpha_{\varepsilon 0} + n_{\varepsilon 0},\; \alpha_{\varepsilon 1}
+ n_{\varepsilon 1})$, where $n_{\varepsilon j}$ counts observations in
$B_{\varepsilon j}$.  This update is exact and closed-form for any finite
truncation of the tree, and yields closed-form marginal likelihoods
\citep{BergerGuglielmi2001, HolmesChoudhuriBhattacharya2003} --- a
property particularly valuable for the model selection and model averaging
computations in Chapter~\ref{chap:metaanalysis}.  In
Chapter~\ref{chap:metaanalysis} a finite-depth truncation is used with
depth chosen to balance approximation error against computational cost,
following the adaptive truncation strategy of
\citet{TrippaHolmesJohnson2011}.  The DP handles the combinatorial
question of which studies share a common data-generating density, while the
P\'{o}lya tree handles the continuous question of what that density looks
like, without committing to a parametric family.

% =============================================================================
\section{The three papers}
\label{sec:intro:papers}
% =============================================================================

% -----------------------------------------------------------------------------
\subsection{Paper~1 --- Learning dynamically: recursive estimation via gradient flows}
\label{subsec:intro:paper1}
% -----------------------------------------------------------------------------

Newton's algorithm \citep{NewtonRaftery1994} is a sequential estimator for
mixture models: given a current estimate $G_n \in \cP(\Th)$ of the mixing
measure, the update
\begin{equation}
  G_{n+1}(d\theta)
  \;=\;
  (1 - \alpha_n)\, G_n(d\theta)
  \;+\;
  \alpha_n\,
  \frac{k(x_{n+1}, \theta)\, G_n(d\theta)}
       {\displaystyle\int k(x_{n+1}, \theta)\, G_n(d\theta)}
  \label{eq:newton-recursion}
\end{equation}
incorporates the new observation $x_{n+1}$ through a likelihood-weighted
reweighting of the current estimate.  The algorithm is computationally
attractive and has been shown to converge almost surely to the true mixing
measure under conditions including correct model specification and a
decreasing step-size sequence satisfying $\sum_n \alpha_n = \infty$ and
$\sum_n \alpha_n^2 < \infty$ \citep{TokdarMartinGhosh2009}.  Yet its
Bayesian interpretation has remained unclear: \citet{FortiniPetrone2020}
establish that no choice of $\alpha_n$ makes $G_n$ equal to the posterior
mean $\bE[G \mid x_{1:n}]$ under any species sampling prior, and the
predictive sequence induced by Newton's recursion does not in general
satisfy the martingale property required for a valid martingale posterior.

Chapter~\ref{chap:gradientflows} resolves this by establishing that
Newton's recursion is the forward Euler discretisation of the Fisher-Rao
gradient flow of the marginal log-likelihood
\begin{equation}
  \cF(G \mid x_t)
  \;=\;
  -\log \int k(x_t, \theta)\, G(d\theta),
  \label{eq:marginal-loglik}
\end{equation}
a functional on $(\cP(\Th), \dH)$ where $\dH$ denotes the Hellinger
metric.  This identification has three consequences.  First, it provides a
decision-theoretic interpretation: Newton minimises a sequential prediction
loss, and the step sizes $\alpha_n$ control the discretisation error, not
the prior strength.  Second, it explains why prior centring fails: the
objective $\cF$ contains no prior term, so no choice of $\alpha_n$ can
make $G_n$ track the posterior mean under a non-trivial prior $\pi$.
Third, it suggests a principled fix: augmenting $\cF$ with a divergence
term $\lambda\, D(G \,\|\, \pi)$ recovers a generalised variational
inference problem \citep{KnoblauchJewsonDamoulas2022} over the
Wasserstein--Fisher-Rao geometry, whose particle system implementation
converges to a better-calibrated estimate and explores the parameter space
more thoroughly at small sample sizes.  The gradient flow language further
unifies Newton's recursion with martingale posteriors and score-based
sequences \citep{WalkerHolmes2024}, which
Chapter~\ref{chap:gradientflows} identifies as gradient flows under
different metrics on $\cP(\Th)$, each corresponding to a different implicit
loss function in the decision-theoretic framework of
Section~\ref{subsec:intro:bayes}.

% -----------------------------------------------------------------------------
\subsection{Paper~2 --- Learning efficiently: coresets for random partitions}
\label{subsec:intro:paper2}
% -----------------------------------------------------------------------------

Coreset construction is the problem of finding a small weighted subsample
of the data such that inference on the coreset approximates inference on
the full dataset.  For parametric exponential family models, coresets can
be constructed by sensitivity sampling --- assigning each observation an
importance weight proportional to its maximum influence on the
log-likelihood --- with approximation guarantees scaling as $O(\log n /
\varepsilon^2)$ in the coreset size \citep{HugginsCampbellBroderick2016,
CampbellBroderick2018}.  For Bayesian nonparametric models with
exchangeable random partition structure, these constructions fail: the
log-likelihood does not decompose as a sum of independent per-observation
terms, because the probability of any observation depends on which cluster
it joins, which in turn depends on all other observations through the EPPF.
Standard sensitivity weights, derived from individual likelihoods, cannot
capture this dependence and produce coreset posteriors whose cluster size
distributions diverge from those of the full-data posterior.

Chapter~\ref{chap:coresets} develops a coreset construction for species
sampling models based on predictive matching.  The criterion is
decision-theoretic: a weighted subsample $(x_{i_1}, w_1), \ldots,
(x_{i_m}, w_m)$ is a valid coreset if the predictive sequence it induces
approximates the predictive sequence of the full data to accuracy
$\varepsilon$ in a suitable divergence, with the coreset size $m$ growing
sublinearly in $n$.  This criterion connects directly to
Section~\ref{subsec:intro:predictive}: a Bayesian learner that has
compressed its data to a coreset makes the same predictions as one that has
retained everything, up to controlled error.  The construction exploits the
exchangeability of the EPPF and its recursive product form, which together
determine how the cluster size distribution must be preserved in the
subsample to maintain predictive accuracy across all $n$ steps.

% -----------------------------------------------------------------------------
\subsection{Paper~3 --- Learning across studies: meta-analysis via DP mixtures of P\'{o}lya trees}
\label{subsec:intro:paper3}
% -----------------------------------------------------------------------------

Meta-analysis is the problem of combining evidence from multiple
independent studies to draw inference about a common underlying
phenomenon.  The canonical parametric approach --- the
DerSimonian--Laird random effects model \citep{DerSimonianLaird1986} ---
assumes that study-level effects are drawn from a Gaussian distribution,
an assumption that is routinely violated in practice: studies differ in
population characteristics, measurement protocols, and outcome definitions
in ways that produce non-Gaussian and often multimodal heterogeneity
\citep{HigginsThompson2002, TurnerEtAl2012}.  Forcing a Gaussian model on
this heterogeneity can produce pooled estimates that do not correspond to
any study's population and confidence intervals that are systematically too
narrow.

Chapter~\ref{chap:metaanalysis} proposes a nonparametric approach based on
Dirichlet process mixtures of P\'{o}lya trees.  Each study $j$ generates
data from an unknown density $f_j \in \cF$, and the collection
$(f_1, f_2, \ldots)$ is modelled as i.i.d.\ draws from an unknown
distribution $\cQ$ over the space of densities $\cF$, with $\cQ \sim
\DP{\alpha}{\cQ_0}$ where $\cQ_0$ is a base measure on $\cF$ given by a
P\'{o}lya tree prior $\PT{\cB}{\cA}$.  The inferential targets are the
cluster assignments of studies, the pooled density within each cluster, and
the between-study heterogeneity as characterised by the distribution $\cQ$.
The DP prior on $\cQ$ handles the combinatorial question of which studies
share a common data-generating density, without requiring the number of
clusters to be specified in advance.  Within each cluster, the P\'{o}lya
tree prior on $f_j$ allows flexible density estimation without parametric
constraints, while conjugacy enables closed-form marginal likelihoods at
each node of the tree.  The model is applied to a collection of $k$
randomised controlled trials from \texttt{[dataset]}, providing
\texttt{[main substantive finding]}.

% =============================================================================
\section{Thesis roadmap}
\label{sec:intro:roadmap}
% =============================================================================

This thesis is organised as follows.
Chapter~\ref{chap:gradientflows} develops the gradient flow theory for
recursive mixture estimation.  The main result is that Newton's algorithm
is the forward Euler discretisation of the Fisher-Rao gradient flow of the
marginal log-likelihood on $(\cP(\Th), \dH)$; the reader should take away
that the gradient flow language provides a decision-theoretic foundation
for recursive estimators and a principled route to incorporating prior
information.
Chapter~\ref{chap:coresets} develops a coreset theory for species sampling
models based on predictive matching.  The main result is a construction
that preserves the EPPF of the full dataset within a weighted subsample of
sublinear size; the reader should take away that Bayesian compression for
partition models requires a predictive criterion rather than a
likelihood-based one.
Chapter~\ref{chap:metaanalysis} develops a DP mixture of P\'{o}lya trees
model for nonparametric meta-analysis and applies it to \texttt{[dataset]}.
The main result is a fully nonparametric pooling procedure that clusters
studies and estimates within-cluster densities without parametric
constraints; the reader should take away that nonparametric priors can
resolve the heterogeneity misspecification problem that affects standard
random effects models.
Chapter~\ref{chap:conclusions} concludes with a discussion of open
problems, including the extension of the gradient flow framework to other
BNP models, the development of adaptive coreset constructions that update
online, and the use of infinite-dimensional priors as regularisers in
large-scale machine learning.  We do not address model selection across the
three frameworks, nor do we provide a unified computational implementation;
these remain open directions.

% =============================================================================
% Bibliography
% =============================================================================
% \printbibliography   % uncomment when .bib file is attached

\end{document}